\chapter{Evaluation and Discussion}

\section{System Performance}

The April 2026 audit prioritised repository-backed evidence over unreproducible
prototype measurements. Consequently, this chapter distinguishes between
observations that remain defensible from the current codebase and earlier
lab-style measurements that should now be treated as provisional until they are
rerun under a documented procedure.

\subsection{End-to-End Command Latency}

End-to-end command latency should be defined as the elapsed time from a user
command in the client interface to the corresponding physical relay transition
at the ESP32 node. Earlier drafts included exact stopwatch-based measurements,
but those values were not re-measured during the April 2026 audit and should
not be presented as current verified benchmarks.

From the implemented architecture, the main latency contributors remain easy to
identify: frontend request dispatch, Flask API processing, MQTT broker
round-trip time, publish-to-subscribe propagation, and physical relay actuation.
This path is still short enough to support practical interactive control in a
prototype setting, but the exact milliseconds must be rerun on real hardware
before any strong quantitative claim is retained.

\subsection{Sensor Data Throughput}

The repository confirms a lightweight telemetry pipeline: the ESP32 firmware is
designed to publish periodic sensor readings, the backend MQTT gateway persists
them through the repository layer, and clients consume updates through REST and
Socket.IO. However, no controlled stress test was rerun during the April 2026
audit, so earlier throughput and database-write figures should be treated as
historical prototype observations rather than as current validated metrics.

\subsection{Alfred AI Response Latency}

Alfred response latency depends on the local Ollama service (or upstream Gemini fallback), network
conditions, prompt size, and whether the request occurs in a warm or cold
session. The April 2026 audit confirmed that Alfred remains integrated into the
current backend and client flows, but it did not rerun a controlled latency
benchmark. Therefore, earlier exact median and percentile figures are better
treated as indicative prototype observations than as present-tense validated
results.

For the current thesis revision, the defensible claim is qualitative: Alfred is
architecturally capable of interactive responses within a human-usable control
workflow, but a fresh benchmark is still required before citing exact response
time statistics.

\subsection{ML Anomaly Detection Inference Time}

The Isolation Forest model inference pipeline (feature assembly, \texttt{decision\_function},
risk classification, alert emission) remains lightweight at the code-structure level:

\begin{itemize}
    \item A lightweight benchmark pipeline exists in the repository and indicates that inference can be executed at interactive speeds on commodity hardware.
    \item The exact latency figures reported in earlier drafts were not re-measured during the April 2026 audit and should therefore be treated as provisional until the benchmark is rerun under a controlled procedure.
    \item End-to-end alert propagation remains architecturally short because inference, persistence, and Socket.IO emission occur in the same backend process.
\end{itemize}

Even without restating an exact millisecond figure, the pipeline remains appropriate
for prototype-level cardiac monitoring because adverse changes in this use case evolve
over seconds to minutes rather than microseconds.

\section{Security Analysis}

The backend addresses the most critical web application security risks as follows:

\begin{description}
    \item[Authentication and access control] Signed bearer-token authentication is
    enforced on protected endpoints through a custom \texttt{@auth\_required}
    decorator. Tokens are generated and verified with \texttt{itsdangerous}, and
    failed login attempts are rate-limited to 10 per minute and 50 per hour.

    \item[Injection prevention] All database access is mediated through SQLAlchemy's
    ORM with parameterised queries. No raw SQL string concatenation is present in
    the codebase, preventing SQL injection (OWASP A03:2021).

    \item[MQTT transport security] The EMQX broker is accessed over TLS on port
    8883. As of the April 2026 hardening pass, the tracked firmware no longer
    contains embedded credentials and supports certificate-based validation via a
    local secrets header, while still permitting an explicit development-only
    insecure fallback during lab setup.

    \item[Rate limiting] Flask-Limiter enforces global limits of 500 requests per
    day and 100 per hour per client IP, with stricter per-endpoint limits for login
    and registration to reduce brute-force abuse.

    \item[Password storage] User passwords are stored as salted hashes generated
    by Werkzeug's password helpers; plaintext passwords are never persisted or logged.
\end{description}

One acknowledged limitation is that the Flask API is served over HTTP (not HTTPS)
within the thesis LAN environment. In a production deployment, TLS termination
via an Nginx reverse proxy with Let's Encrypt certificates would be required.

\section{Usability Assessment}

\subsection{Study Design}

A small formative usability exercise informed the interface iteration of this
prototype. The original participant notes covered both caregiver-like users and
older adults, and the tasks focused on finding live sensor data, toggling a
device, querying Alfred, and navigating to the mobile health view. These notes
remain useful for design reflection, but they do not constitute a rigorous,
reproducible user study by the standards of a formal HCI evaluation.

The original task framing was:

\begin{enumerate}
    \item View the current room temperature reading.
    \item Toggle the living-room light ON via the web dashboard.
    \item Ask Alfred about the current heart-rate reading.
    \item Navigate to the health history screen on the mobile application.
\end{enumerate}

Task completion and time-on-task were observed informally, and participants gave
lightweight subjective feedback after completion. Because the underlying raw
records were not reproduced in the April 2026 audit, the current thesis treats
the usability evidence as qualitative guidance rather than as a source of exact
statistical claims.

\subsection{Task Completion and Rating Results}

The most reliable conclusion from the formative notes is directional rather than
numerical. Simple observation and control tasks appeared easier for users than
chat-oriented or navigation-heavy tasks. That pattern is consistent with the
design changes implemented later in the repository, especially the addition of
voice input for Alfred and the clearer separation of major mobile screens.

\subsection{Qualitative Observations}

The Alfred chat task appeared to be the least intuitive activity for older users
who were unfamiliar with a text-chat interaction style. Informal notes suggest
that several users expected a more direct voice-driven interaction rather than a
typed message workflow.

The health history navigation task also exposed discoverability issues around tab
labels and screen transitions. These observations support the later addition of
voice input for Alfred and the emphasis on clearer screen segmentation in the
current mobile interface.

Caregiver-oriented users adapted to the dashboard metaphor more quickly than the
elderly participants. This asymmetry reinforces a central design lesson of the
project: clinical or caregiver power-user features can coexist with elderly-care
support only if accessibility accommodations are treated as first-class design
requirements rather than as cosmetic additions.

\subsection{Comparative Summary}

Table~\ref{tab:compare} extends the literature comparison from Chapter~2 to include
implementation characteristics and the present validation focus of this work.

\begin{table}[ht]
\centering
\caption{Extended comparison with related systems}
\label{tab:compare}
\begin{tabular}{lccc}
\toprule
\textbf{Criterion} & \textbf{Ours} & \cite{Alam2012} & \cite{Suryadevara2014} \\
\midrule
Open-source        & Yes    & Yes    & No \\
Protocol           & MQTT   & ZigBee & Wi-Fi \\
ML anomaly detect. & IF     & None   & Fuzzy \\
AI assistant       & Alfred (Ollama/Gemini) & None   & None \\
Map/UI editor      & Yes    & No     & No \\
Mobile app         & Flutter & None  & None \\
Current validation focus & Build/test evidence & N/R & N/R \\
\bottomrule
\multicolumn{4}{l}{\footnotesize IF = Isolation Forest;\quad N/R = Not reported}
\end{tabular}
\end{table}

This system differentiates itself primarily through the combination of a generative AI
assistant, an interactive map editor, and a fully open-source end-to-end
implementation. The use of MQTT over TLS with an EMQX Cloud broker provides
deployment flexibility beyond local-Wi-Fi-only solutions, at the cost of requiring an
internet connection for broker communication.

\subsection{Comparison with Previous Studies}

Comparing the present work against the closest prior systems reveals both points of
continuity and substantive advances. \textcite{Alam2012} demonstrated the feasibility
of low-power sensor networks for elderly monitoring using ZigBee mesh radio, reporting
successful activity detection across multiple rooms; however, their architecture did not
include any anomaly-detection model, caregiver alert pipeline, or AI-assisted
interaction layer. \textcite{Suryadevara2014} moved further by applying fuzzy-logic
inference to infer wellness states from behavioural cues (appliance usage patterns),
achieving rule-based classification without requiring wearable sensors. Neither work
integrated continuous physiological monitoring: both systems relied solely on
environmental or behavioural cues and did not monitor the user's heart rate, heart-rate
variability, or any other direct physiological signal. Neither work published a
quantitative test suite or an automated regression strategy.

The present system advances the state of the art along four measurable dimensions.
First, it introduces continuous wearable physiological monitoring via the Coospo H6 BLE
chest strap, filling the most significant clinical gap in both reference systems: the
inability to detect adverse cardiac events that leave no behavioural or environmental
trace. Second, it replaces hand-crafted rule sets with a trained Isolation Forest model
that classifies physiological risk across four severity levels (normal, caution,
warning, critical) from real-time heart-rate and HRV signals. Third, it integrates a
generative AI assistant (Alfred) that enables natural-language device control in both
Vietnamese and English, a feature neither prior system considered. Fourth, the
implementation is accompanied by all 448 backend unit and integration tests passing
(75\% line coverage achieved) and a 205-case Flutter regression suite of which 199 pass,
providing a degree of reproducibility evidence that the earlier works, constrained by
the tools of their era, did not report.

Where this work falls short of those studies is in longitudinal field evidence: both
\textcite{Alam2012} and \textcite{Suryadevara2014} included multi-day or multi-week
deployment observations in real homes, whereas the present prototype validation remains
laboratory-scale. Closing that gap is the primary direction identified for future work.

\section{Patient Data Privacy and Protection}
\label{sec:privacy}

\subsection{Sensitivity of Stored Data}

The system stores and processes several categories of personal health data that
warrant specific privacy considerations:

\begin{itemize}
    \item \textbf{Physiological data}: Heart-rate readings (BPM), heart-rate variability
          metrics (RMSSD, SDNN, pNN50), risk classifications, and anomaly alerts stored
          in \texttt{patient\_hr\_records} and \texttt{alerts}.
    \item \textbf{Profile data}: Patient full name, age, gender, and clinical notes
          stored in \texttt{patient\_profiles}.
    \item \textbf{Behavioural data}: Motion sensor events (PIR), device control history
          (\texttt{control\_log}), and room occupancy patterns derivable from the floor-plan map.
    \item \textbf{Account credentials}: Email addresses and hashed passwords in
          \texttt{users}.
\end{itemize}

In jurisdictions where health data is regulated (e.g., under Vietnam's Cybersecurity Law
2018, the EU General Data Protection Regulation, or equivalent frameworks), several of
these categories qualify as sensitive personal data requiring explicit legal basis for
processing.

\subsection{Current Protection Measures}

The prototype implements the following safeguards:

\begin{description}
    \item[Authentication and access control] All health data endpoints require a
          valid bearer token. Queries are scoped to the authenticated user's
          \texttt{user\_id}, preventing cross-user data access. The admin panel
          additionally requires the \texttt{admin} role to access aggregate user lists.

    \item[Password protection] Passwords are stored as salted hashes using
          Werkzeug's \texttt{generate\_password\_hash()} using its default hashing scheme (scrypt in Werkzeug $\geq$ 3.0).
          Plaintext passwords are never written to disk or logged.

    \item[Transport security] The REST API is served over HTTPS in the Railway
          cloud deployment (TLS terminated at the platform edge). The MQTT broker
          connection uses TLS on port 8883. The LAN prototype uses HTTP
          between the Vite dev server and backend, which is an acknowledged
          limitation for production use.

    \item[Account deletion] The \texttt{DELETE /api/auth/account} endpoint implements
          self-service account cancellation using a soft-delete approach: the
          account is deactivated (\texttt{is\_active=False}) and personal
          identifiers (username, email, password hash) are anonymised, preventing
          further login while retaining referential integrity of associated
          records. This provides a functional ``right to erasure'' workflow
          consistent with data minimisation principles.

    \item[Audit trail] The \texttt{control\_log} table retains a tamper-evident
          record of every device command with its origin source, allowing
          post-incident review without exposing raw health records.
\end{description}

\subsection{Acknowledged Gaps and Future Requirements}

The following protections are absent in the current prototype and represent
requirements for a production-grade deployment:

\begin{enumerate}
    \item \textbf{Encryption at rest}: The MySQL database is not encrypted at the
          storage layer. A production deployment should enable MySQL Transparent
          Data Encryption (TDE) or an equivalent mechanism.

    \item \textbf{Data retention policy}: No automated deletion schedule is implemented
          for aged sensor readings or resolved alerts. A production system should
          define and enforce a retention window (e.g., delete sensor readings older
          than 12 months) to limit unnecessary accumulation of sensitive data.

    \item \textbf{Formal consent workflow}: The current registration flow collects
          only a username, email, and password. A compliant deployment would
          require explicit informed consent for health data processing before the
          first sensor reading is stored.

    \item \textbf{Data export}: There is no mechanism for a user to export their
          complete health record in a portable format (e.g., JSON or CSV). A right
          to data portability would require implementing a \texttt{GET /api/auth/export}
          endpoint.

    \item \textbf{Security audit logging}: User login events, failed authentication
          attempts, and administrative actions (role changes, account deletions)
          are not currently written to a structured security log. A production
          deployment should centralise these events for intrusion detection and
          compliance reporting.
\end{enumerate}

These gaps are consistent with the scope of a thesis prototype and are noted
explicitly so that any future deployment proceeds with a full Data Protection
Impact Assessment (DPIA) rather than treating the prototype security posture as
production-ready.

\section{Summary of Evaluation}

Table~\ref{tab:eval_summary} provides a consolidated summary of all evaluation
dimensions.

\begin{table}[ht]
\centering
\caption{Consolidated evaluation summary}
\label{tab:eval_summary}
\begin{tabular}{lll}
\toprule
\textbf{Dimension} & \textbf{Target} & \textbf{Achieved} \\
\midrule
Backend validation           & Clean startup and targeted tests & Repository checks passed \\
Web validation               & Production build integrity & Vite build passed \\
Mobile validation            & Static analysis and regression tests & \texttt{flutter analyze} passed; 199 of 205 tests pass \\
Auth and security model      & Consistent implementation & Hardened and documentation-aligned \\
Firmware compile status      & Fresh rebuild in current audit & Pending toolchain reproduction \\
Hardware latency benchmark   & Reproducible physical measurement & Pending re-measurement \\
AI latency benchmark         & Controlled benchmark rerun & Pending re-measurement \\
Field reliability study      & Long-duration deployment evidence & Pending future work \\
\bottomrule
\end{tabular}
\end{table}
